\documentclass[12pt,a4paper]{article}
\usepackage{amsmath,amssymb,mathrsfs,tikz,times,pifont}
\usepackage{enumitem}
\newcommand\circitem[1]{%
\tikz[baseline=(char.base)]{
\node[circle,draw=gray, fill=red!55,
minimum size=1.2em,inner sep=0] (char) {#1};}}
\newcommand\boxitem[1]{%
\tikz[baseline=(char.base)]{
\node[fill=cyan,
minimum size=1.2em,inner sep=0] (char) {#1};}}
\setlist[enumerate,1]{label=\protect\circitem{\arabic*}}
\setlist[enumerate,2]{label=\protect\boxitem{\alph*}}
%%%::::::by chnini ameur :::::::%%%
\everymath{\displaystyle}
\usepackage[left=1cm,right=1cm,top=1cm,bottom=1.7cm]{geometry}
\usepackage{array,multirow}
\usepackage[most]{tcolorbox}
\usepackage{varwidth}
\tcbuselibrary{skins,hooks}
\usetikzlibrary{patterns}
%%%::::::by chnini ameur :::::::%%%
\newtcolorbox{exa}[2][]{enhanced,breakable,before skip=2mm,after skip=5mm,
colback=yellow!20!white,colframe=black!20!blue,boxrule=0.5mm,
attach boxed title to top left ={xshift=0.6cm,yshift*=1mm-\tcboxedtitleheight},
fonttitle=\bfseries,
title={#2},#1,
% varwidth boxed title*=-3cm,
boxed title style={frame code={
\path[fill=tcbcolback!30!black]
([yshift=-1mm,xshift=-1mm]frame.north west)
arc[start angle=0,end angle=180,radius=1mm]
([yshift=-1mm,xshift=1mm]frame.north east)
arc[start angle=180,end angle=0,radius=1mm];
\path[left color=tcbcolback!60!black,right color = tcbcolback!60!black,
middle color = tcbcolback!80!black]
([xshift=-2mm]frame.north west) -- ([xshift=2mm]frame.north east)
[rounded corners=1mm]-- ([xshift=1mm,yshift=-1mm]frame.north east)
-- (frame.south east) -- (frame.south west)
-- ([xshift=-1mm,yshift=-1mm]frame.north west)
[sharp corners]-- cycle;
},interior engine=empty,
},interior style={top color=yellow!5}}
%%%%%%%%%%%%%%%%%%%%%%%

\usepackage{fancyhdr}
\usepackage{eso-pic}         % Pour ajouter des éléments en arrière-plan
% Commande pour ajouter du texte en arrière-plan
\AddToShipoutPicture{
    \AtTextCenter{%
        \makebox[0pt]{\rotatebox{80}{\textcolor[gray]{0.6}{\fontsize{5cm}{5cm}\selectfont PGB}}}
    }
}
\usepackage{lastpage}
\fancyhf{}
\pagestyle{fancy}
\renewcommand{\footrulewidth}{1pt}
\renewcommand{\headrulewidth}{0pt}
\renewcommand{\footruleskip}{10pt}
\fancyfoot[R]{
\color{blue}\ding{45}\ \textbf{2026}
}
\fancyfoot[L]{
\color{blue}\ding{45}\ \textbf{Prof:M. BA}
}
\cfoot{\bf
\thepage /
\pageref{LastPage}}
\begin{document}
\renewcommand{\arraystretch}{1.5}
\renewcommand{\arrayrulewidth}{1.2pt}
\begin{tikzpicture}[overlay,remember picture]
\node[draw=blue,line width=1.2pt,fill=purple,text=blue,inner sep=3mm,rounded corners,pattern=dots]at ([yshift=-2.5cm]current page.north) {\begingroup\setlength{\fboxsep}{0pt}\colorbox{white}{\begin{tabular}{|*1{>{\centering \arraybackslash}p{0.28\textwidth}} |*2{>{\centering \arraybackslash}p{0.2\textwidth}|} *1{>{\centering \arraybackslash}p{0.19\textwidth}|} }
\hline
\multicolumn{3}{|c|}{$\diamond$$\diamond$$\diamond$\ \textbf{Lycée de Dindéfélo}\ $\diamond$$\diamond$$\diamond$ }& \textbf{A.S. : 2025/2026} \\ \hline
\textbf{Matière: Mathématiques}& \textbf{Niveau : T}\textbf{S2} &\textbf{Date: 18/05/2025} & \textbf{Durée : 4 heures} \\ \hline
\multicolumn{4}{|c|}{\parbox[c]{10cm}{\begin{center}
\textbf{{\Large\sffamily Devoir n$ ^{\circ} $ 1 Du 2$ ^\text{\bf nd} $ Semestre}}
\end{center}}} \\ \hline
\end{tabular}}\endgroup};
\end{tikzpicture}
\vspace{3cm}

\section*{\underline{Exercice 1 : Probabilités} (5,75 points) $\approx$ 36 mns)}

Le mois d'Août est très pluvieux. Le premier Août 2020, il a plu. Une étude a montré que :
\begin{itemize}
    \item S'il a plu un jour, la probabilité qu'il pleuve le jour suivant est $1/3$.
    \item S'il n'a pas plu un jour, la probabilité qu'il ne pleuve pas le jour suivant est de $2/5$.
\end{itemize}

On note par :
\begin{itemize}
    \item $A_n$ : "Il a plu au $n$-ième jour du mois d'Août"
    \item et $P_n = P(A_n)$
\end{itemize}

\begin{enumerate}
    \item Donner $P_1$. \hfill \textbf{(0,25 pt)}
    \item Calculer $P_2$. \hfill \textbf{(0,75 pt)}
    \item Montrer que $P_{n+1} = -\frac{4}{15} P_n + \frac{3}{5}$. \hfill \textbf{(1,5 pt)}
    \item On note $u_n = P_n - \frac{9}{19}$.
    \begin{enumerate}
        \item Montrer que $(u_n)$ est une suite géométrique. \hfill \textbf{(1 pt)}
        \item Exprimer $P_n$ en fonction de $n$. \hfill \textbf{(1,25 pt)}
    \end{enumerate}
    \item Quelle est la probabilité qu'il pleuve le 5 Août ? \hfill \textbf{(1 pt)}
\end{enumerate}

\section*{\underline{Exercice 2 :} (04 points $\approx$ 30 mns)}
\noindent On considère la suite numérique $(u_n)_{n \in \mathbb{N}}$ définie par : 
$\begin{cases} 
U_0 = 6 \\ 
U_{n+1} = \dfrac{1}{U_n} + \dfrac{3}{4}U_n, \, n \in \mathbb{N} 
\end{cases}$

\begin{enumerate}
    \item Calculer $u_1$ et $u_2$. \hfill \textbf{(0,5 pt)}
    \item Démontrer par récurrence que : $\forall n \in \mathbb{N}, \quad u_n \ge \sqrt{3}$. \hfill \textbf{(1 pt)}
    \item Soit $f$ la fonction définie sur $]0, +\infty[$ par $f(x) = \dfrac{1}{x} + \dfrac{3}{4}x$.
    \begin{enumerate}
        \item Étudier le sens de variations de $f$. \hfill \textbf{(1 pt)}
        \item En déduire par récurrence que $(u_n)_{n \in \mathbb{N}}$ est strictement décroissante. \hfill \textbf{(0,5 pt)}
    \end{enumerate}
    \item Montrer que $(u_n)_{n \in \mathbb{N}}$ est convergente et déterminer sa limite. \hfill \textbf{(1 pt)}
\end{enumerate}
\section*{\underline{Problème :} ( 10,25 points $\approx$ 144 mns)}

\textbf{Partie A : } Soit les fonctions $g$ et $h$ définies respectivement par : $g(x) = x^2 + x - 2 + 2\ln(x)$ sur $]0; +\infty[$ et $h(x) = 1 - (x+1)^2e^x$ sur $]-\infty; 0]$.

\begin{enumerate}
    \item Étudier les variations de $g$ et $h$ respectivement sur leurs domaines de définitions. \hfill \textbf{(1,5 pt)}
    \item Calculer $g(1)$ puis déduire le signe de $g$ sur $]0; +\infty[$. \hfill \textbf{(0,5 pt)}
    \item Donner le signe de $h$ sur $]-\infty; 0]$. \hfill \textbf{(0,5 pt)}
\end{enumerate}

\vspace{0.5cm}

\textbf{Partie B : } Soit $f$ la fonction définie par : 
$f(x) = 
\begin{cases} 
x + \left(1 - \dfrac{2}{x}\right) \ln(x) & \text{si } x > 0 \\
x + 1 - (x^2 + 1)e^x & \text{si } x \le 0 
\end{cases}$

\begin{enumerate}
    \item Montrer que le domaine de définition $D_f = \mathbb{R}$. Calculer les limites aux bornes de $D_f$. \hfill \textbf{(0,75 pt)}
    \item Étudier les branches infinies. \hfill \textbf{(1 pt)}
    \item Étudier la position de $(\mathcal{C}_f)$ par rapport à $(\Delta)$ d'équation $y = x + 1$ en $-\infty$. \hfill \textbf{(0,5 pt)}
    \item Étudier la continuité de $f$ en 0. \hfill \textbf{(0,25 pt)}
    \item $f$ est-elle dérivable à droite de 0 ? Étudier la dérivabilité à gauche de 0 puis interpréter le résultat. \hfill \textbf{(0,75 pt)}
    \item Montrer que $\forall x > 0, f'(x) = \dfrac{g(x)}{x^2}$. En déduire le signe de $f'(x)$ sur $]0; +\infty[$. \hfill \textbf{(1 pt)}
    \item Montrer que $\forall x < 0, f'(x) = h(x)$. En déduire le signe de $f'(x)$ sur $]-\infty; 0[$. \hfill \textbf{(1 pt)}
    \item Dresser le tableau de variation de $f$ sur $D_f$. Puis tracer $(\mathcal{C}_f)$. \hfill \textbf{(1,5 pt)}
\end{enumerate}

\vspace{0.5cm}

\textbf{Partie C : } Soit $k$ la restriction de $f$ sur $I = ]1; +\infty[$.

\begin{enumerate}
    \item Montrer que $k$ réalise une bijection de $I$ vers un intervalle $J$ à préciser. \hfill \textbf{(0,25 pt)}
    \item $k^{-1}$ est-elle dérivable en 1 ? $k^{-1}$ est-elle dérivable sur $J$ ? \hfill \textbf{(0,5 pt)}
    \item Tracer la courbe de $(\mathcal{C}_{k^{-1}})$ \hfill \textbf{(0,25 pt)}
\end{enumerate}
\end{document}
